% ============================================================================
% CH 3 — RELATED WORK                                    [slug: related-work]
% STATUS: reading list largely covered (S2 survey re-read in progress); grows
% with parallel literature sessions. Argumentative gap analysis, NOT a survey:
% every subsection states how the works relate to ours; chapter ENDS with the gap.
% Material map: ../outline.md#related-work
% ============================================================================
\cleardoublepage
\chapter{Related Work (7--9)}   % page goal — scaffolding, DELETE before submission
\label{ch:related}

\section{Scope: Two Axes}
\label{sec:rw-scope}
% What was considered and why: axis 1 = feature-aware backbone (how features
% enter MF), axis 2 = explanation machinery (how models earn interpretability).

\section{Feature-Aware Matrix Factorization}
\label{sec:rw-feature-mf}
% CMF, RLFM, SVDFeature, FM family (AFM/NFM/DeepFM/Wide&Deep), CTR/CDL,
% DropoutNet; "expressive but opaque" foil; grounding-bridge verdict (FM/GMF
% not THE mechanism); feature-injection fork: prediction-side vs prototype-side.

\section{Prototype-Based Interpretable Models}
\label{sec:rw-prototypes}
% ProtoPNet lineage, ProtoTree, part-prototype challenges (S5); PPM prominence
% as credibility anchor; Axis-A/B mapping and disclaimers (2026-06-16_2057).

\section{Prototypes in Recommendation}
\label{sec:rw-proto-recsys}
% ACF (conceptual predecessor), ProtoCF (closest prior art), UIPC-MF (closest
% competitor), ProtoMF forward citations (I1-I5): none grounds prototypes in
% catalogue features.

\section{Concept Grounding and Semantic Anchoring}
\label{sec:rw-concepts}
% CBMs, SENN, attention-as-explanation debate (E4a/E4b); dc02's enforced
% grounding as the contrast concept.

\section{Evaluating Explanations and Interpretability}
\label{sec:rw-eval-explanations}
% S4 measuring-why, G1 fidelity metrics; feeds Ch 10's metric choices.

\section{The Gap}
\label{sec:rw-gap}
% Closing move (DBIS/VUB norm): no ProtoMF descendant grounds prototypes in
% item features intrinsically; modality-light positioning (catalogue attributes
% only, vs review-corpus methods); this thesis fills it with a staged lineage.
